\section{Aprendizado de Máquina para Detecção de Anomalias}
\label{sec:ml-conceitos}

O uso de aprendizado de máquina é coerente com o estado da arte recente.
Trabalhos como \cite{farhad2023} demonstram que o aprendizado de máquina em
LoRaWAN vem sendo explorado para otimizar parâmetros de rede
(\textit{spreading factor}, potência de transmissão) e classificar padrões
de mobilidade. Embora parte significativa da literatura esteja concentrada em
desempenho e eficiência, esse mesmo conjunto de atributos também é relevante
para a detecção de anomalias, já que comportamentos suspeitos se refletem em
mudanças desses padrões operacionais \cite{farhad2023, vangelista2023}.

Esta seção trata dos conceitos de séries temporais e de detecção de anomalias
usados pela frente de aprendizado de máquina. Eles sustentam o método da
Seção~\ref{subsec:ml-dataset} e os resultados da Seção~\ref{sec:resultados_ml}.

%%%%%%%%%%%%%%%%%%%%%%%%%%%%%%%%%%%%%%%%%%%%%%%%%%%%

\subsection{Por que detectar desvios do RSSI}
\label{subsec:ml-motivacao}

A autenticação de um quadro, pelo código de integridade do LoRaWAN ou pela
\textit{tag} do Ascon usada no protótipo da equipe, mostra que o remetente
conhece a chave da sessão. Ela não mostra de onde o sinal partiu. Se a chave de
um nó sensor for obtida por um atacante, os quadros que ele enviar serão
aceitos como legítimos, e o conteúdo não permite distingui-los dos do nó
verdadeiro. Esse é o caso de \textit{spoofing} que a criptografia não cobre.

O que ainda diferencia o atacante é a camada física. Um transmissor em outra
posição, ou com outra potência, chega ao \textit{gateway} com outro RSSI, e o
\textit{gateway} já mede o RSSI de todo quadro, sem mudança no protocolo. Na
literatura, desvios do RSSI fora da faixa esperada para o nó são tratados como
indício de \textit{spoofing} \cite{pourghasem2025spoofing}, e desvios nos
metadados de rádio, entre eles o RSSI, expõem o \textit{jamming} reativo
\cite{ruotsalainen2022jamming}.

O RSSI, porém, também muda sem ataque, com o canal de frequência, a taxa de
dados e o ambiente (Seção~\ref{subsec:ml-rssi}). Um desvio só tem significado
em relação ao valor esperado para aquele nó naquele momento. Por isso a
detecção deste trabalho prevê o RSSI de cada quadro e alarma quando o valor
observado se afasta da previsão além de um limiar
(Seção~\ref{subsec:ml-deteccao}). Prever, em vez de comparar com um valor
fixo, acompanha as mudanças de nível que acontecem sem ataque, e não exige
exemplos rotulados de ataque para o treino. Como a previsão usa só quadros
passados, a decisão pode ser tomada no próprio \textit{gateway}, quando cada
quadro chega \cite{blazquez2021review}, o que pede um modelo leve o bastante
para um microcontrolador.

\subsection{O RSSI como série temporal}
\label{subsec:ml-rssi}

O RSSI é a potência com que o \textit{gateway} recebe um quadro, em dBm. O
valor é negativo e é mais forte quanto mais perto de zero \cite{ttnrssi}. O nó
sensor não escolhe o \textit{gateway}. Todos os \textit{gateways} ao alcance
recebem o mesmo quadro, e cada um mede o seu próprio RSSI
\cite{brocaarintegration}. Para formar uma série temporal é preciso, portanto,
fixar um par de nó sensor e \textit{gateway}. Como o RSSI é medido no
\textit{gateway}, ele fica fora da proteção dada pela cifragem do
\textit{payload}.

O RSSI muda também sem ataque. O nó sensor troca de canal de frequência de
forma pseudoaleatória a cada transmissão \cite{lorawan103}, e a taxa de dados
adaptativa (ADR) pode mudar a taxa de dados e a potência de transmissão
\cite{ttnadr}. Além disso, um contador de quadros (FCnt) repetido indica o
mesmo quadro enviado de novo, e não uma medida nova \cite{lorawan103}.

\subsection{Autocorrelação e estacionariedade}
\label{subsec:ml-autocorrelacao}

A função de autocorrelação (ACF) mede a relação linear entre a série e ela
mesma deslocada de $k$ passos, chamados de defasagem. A função de
autocorrelação parcial (PACF) mede a mesma relação depois de remover o efeito
das defasagens intermediárias \cite{hyndman2021}.

Uma série é estacionária quando as suas propriedades estatísticas não dependem
do instante observado. A diferenciação, que troca cada valor pela diferença em
relação ao anterior, remove mudanças de nível. No teste KPSS, a hipótese nula é
que a série é estacionária, e um valor-p pequeno indica que é preciso
diferenciar \cite{hyndman2021}.

\subsection{Modelos de previsão}
\label{subsec:ml-modelos}

A referência mais simples é o \textit{random walk}, cuja previsão é a última
observação \cite{hyndman2021}, como na Equação~\ref{eq:ml-random-walk}.
\begin{equation}
	\hat{y}_t = y_{t-1}
	\label{eq:ml-random-walk}
\end{equation}

O ARIMA$(p,d,q)$ combina $p$ termos autorregressivos, $d$ diferenciações e $q$
termos de médias móveis dos erros anteriores \cite{box2015, hyndman2021}. Com o
operador de atraso $B$, tal que $B y_t = y_{t-1}$, o modelo se escreve como na
Equação~\ref{eq:ml-arima}, em que $\varepsilon_t$ é um ruído branco.
\begin{equation}
	\left(1 - \phi_1 B - \dots - \phi_p B^p\right)(1-B)^d y_t = c + \left(1 + \theta_1 B + \dots + \theta_q B^q\right)\varepsilon_t
	\label{eq:ml-arima}
\end{equation}

O \textit{random walk} é o caso ARIMA$(0,1,0)$, e o ARIMA$(0,1,1)$ equivale à
suavização exponencial simples \cite{hyndman2021}. Os coeficientes podem ser
estimados uma vez e depois aplicados, fixos, a uma série mais longa
\cite{statsmodelsapply}. Assim se obtém a previsão um passo à frente, em que
cada valor é previsto com as observações reais até o instante anterior.

\subsection{Avaliação da previsão}
\label{subsec:ml-avaliacao}

Os dados de treino estimam os parâmetros, e os de teste, que não entram no
ajuste, medem a acurácia. Em séries temporais a divisão é cronológica, para que
nenhuma informação futura entre no treino \cite{hyndman2021, bergmeir2012}.
Entre treino e teste, um conjunto de validação serve para as escolhas, como a
ordem do modelo, de modo que o teste seja usado só para a medida final.

O erro de previsão é medido por duas métricas na unidade dos dados, a raiz do
erro quadrático médio (RMSE) e o erro absoluto médio (MAE), definidas nas
Equações~\ref{eq:ml-rmse} e~\ref{eq:ml-mae} \cite{hyndman2021}.
\begin{equation}
	\mathrm{RMSE} = \sqrt{\frac{1}{N}\sum_{t=1}^{N}\left(y_t - \hat{y}_t\right)^2}
	\label{eq:ml-rmse}
\end{equation}
\begin{equation}
	\mathrm{MAE} = \frac{1}{N}\sum_{t=1}^{N}\left|y_t - \hat{y}_t\right|
	\label{eq:ml-mae}
\end{equation}
O RMSE dá mais peso aos erros grandes, e o MAE pesa todos os erros igualmente.
Um método que minimiza o MAE leva a previsões da mediana, e um que minimiza o
RMSE, a previsões da média \cite{hyndman2021}. Uma diferença de ordem entre as
duas métricas indica que a comparação depende do tamanho dos erros
considerados.

\subsection{Detecção de anomalias por previsão}
\label{subsec:ml-deteccao}

Em operação real, dados rotulados que identifiquem quadros associados a
ataques são raros, o que dificulta o uso de classificadores supervisionados.
Por isso a detecção é tratada como um problema de previsão. Um modelo estima o
valor esperado de cada observação a partir das anteriores, e o ponto é
declarado anômalo quando a distância entre o valor observado $x_t$ e o valor
esperado $\hat{x}_t$ passa de um limiar $\tau$ \cite{blazquez2021review,
chalapathy2019}, como na Equação~\ref{eq:ml-deteccao}. Quando $\hat{x}_t$ vem
só de observações passadas, a decisão é tomada sobre cada dado no momento em
que ele chega, o que permite executá-la na borda da rede.
\begin{equation}
	\left| x_t - \hat{x}_t \right| > \tau
	\label{eq:ml-deteccao}
\end{equation}

Sob resíduos normais e sem correlação, $\tau$ pode ser $c$ vezes o desvio
padrão dos resíduos, com $c = 2{,}58$ para 99\% \cite{hyndman2021}. Quando a
suposição de normalidade falha, um limiar não paramétrico, tirado da própria
distribuição dos erros, é uma alternativa \cite{hundman2018spacecraft}. Em
LoRaWAN, o \textit{jamming} reativo aumenta o RSSI e é detectado comparando os
metadados recebidos com os esperados para o nó sensor
\cite{ruotsalainen2022jamming}, e desvios do RSSI fora da faixa esperada são
tratados como indício de \textit{spoofing} \cite{pourghasem2025spoofing}.

\subsection{Trabalhos relacionados e lacuna}
\label{subsec:ml-lacuna}

A literatura sobre o RSSI em redes LoRaWAN se divide em duas linhas que pouco se
encontram. A primeira prevê o RSSI. \citeonline{guerra2024forecasting}
compararam ARIMA, inteligência artificial e modelos híbridos em médias horárias
de oito nós, e o ARIMA com a temperatura como variável externa foi a
alternativa mais competitiva. \citeonline{fhima2026harbors} previram vários
passos à frente o RSSI de nós móveis num porto, reajustando o modelo a cada
ponto. Nenhum dos dois usa a previsão para detectar anomalias, nem trata o
canal de frequência, que \citeonline{harinda2022multigateway} não viram afetar o
RSSI em ambiente interno.

A segunda linha detecta ataques pelos metadados de rádio.
\citeonline{ruotsalainen2022jamming} compara os metadados recebidos com os
esperados para o nó, e \citeonline{pourghasem2025spoofing} classificam quadros
com modelos supervisionados, o que exige rótulos de ataque. Falta, entre as duas
linhas, um modelo leve que forneça o valor esperado de cada quadro, com uma
sensibilidade medida em dB e um alarme falso controlado, e que rode no próprio
\textit{gateway}. Este trabalho ocupa essa lacuna com os metadados de rádio do
\textit{dataset} Tour Perret \cite{tourperret_dataset} e com a execução do
detector num microcontrolador ESP32.
